
Or d'après VI 4.8, $K^{\bullet}(X)$ est la $K^{\bullet}(S)$-algèbre
engendrée par les $\xi_i-1$, soumis aux relations déduites de :
\[
\gamma_t(E-n)=\prod_{i=1}^n\gamma_t(\xi_i-1).
\]
Donc $K^{\bullet}(X)/I=K^{\bullet}(S)/J$, où $J$ est l'idéal engendré par les $\gamma^j(E-n)$ pour $j=1,\ldots,n$, donc le $\lambda$-idéal engendré par $E-n$. Par suite,
$K_\bullet(Y)\simeq K_\bullet(S)/J K_\bullet(S)$. On achève la démonstration par le lemme :

\textbf{Lemme 4.2.} \uline{Soient} $S$ \uline{un schéma noethérien}, $B$ \uline{un sous-groupe parabolique d'un} $S$\uline{-groupe réductif}, $P$ \uline{un torseur sous le radical unipotent} $U$ \uline{de} $B$; \uline{alors on a un isomorphisme} :
\[
K_\bullet(S)\xrightarrow{\sim}K_\bullet(P).
\]
En effet, il existe une suite de sous-schémas en groupes de $U$ :
\[
U_0=U\supset U_1\supset U_2\supset\cdots\supset U_p=e,
\]
fermés dans $B$, finie car $B$ est noethérien, et tels que, pour tout $i$, $U_i/U_{i+1}$ soit isomorphe à un $S$-groupe vectoriel (SGA 3 XXVI 2.1). Par suite, $P\to S$ se factorise par une suite de torseurs sous des groupes vectoriels, d'où le résultat d'après 3.4.

\textbf{Corollaire 4.3.} \uline{Soient} $S$ \uline{un schéma noethérien}, $G=\mathrm{Gl}(n)_S$. \uline{Alors l'homomorphisme canonique} $K_\bullet(S)\longrightarrow K_\bullet(G)$ \uline{est un isomorphisme}.

Résulte immédiatement de 4.1, appliqué à $P=G$, puisqu'alors $\underline{E}=\underline{O}_S^n$.

